%! Inverse Functions — Unit 5, Lesson 5.7 — Exit Ticket
\documentclass[10pt]{article}
\usepackage{algebra2-article}
\usepackage{algebra2-boxes}
\newcommand{\TallMath}[1]{$\displaystyle #1\rule[-1.4em]{0pt}{3.2em}$}

\begin{document}

\pageheader{Unit 5, Lesson 5.7}{Exit Ticket}
\namedateperiod

\vspace{0.08in}

No notes. Reverse the operations \emph{and} the order --- and test anything you build.

\begin{enumerate}[leftmargin=1.8em, itemsep=10pt]
  \item \textbf{Test the pair.} Let $f(x)=4x-8$ and $g(x)=\dfrac{x+8}{4}$.
        \vspace{0.62in}

        (a) $f\big(g(x)\big)=\blank{2.2cm}$ \hspace{0.35in}
        (b) $g\big(f(x)\big)=\blank{2.2cm}$ \hspace{0.35in}
        (c) Inverses? \blank{1.6cm}

        \vspace{4pt}
        (d) Why does the test require \emph{both} orders instead of just one? \blank{5.4cm}

  \item \textbf{Build one, and mind the domain.} Let $f(x)=\sqrt{x-1}$.
        \vspace{0.55in}

        (a) $f^{-1}(x)=\blank{3.0cm}$ \hspace{0.3in}
        (b) Domain of $f^{-1}$: \blank{2.8cm}

        \vspace{5pt}
        (c) In one sentence, say where that restriction comes from --- your sentence should mention
        the \emph{range} of $f$. \writeline

  \item \textbf{SOL-style multiple choice.} If $f(x)=3x+12$, which function is $f^{-1}(x)$?
        \begin{enumerate}[label=\Alph*., leftmargin=1.9em, itemsep=1pt]
          \item $\dfrac{x-12}{3}$
          \item $\dfrac{x}{3}-12$
          \item $\dfrac{1}{3x+12}$
          \item $3x-12$
        \end{enumerate}
\end{enumerate}

\end{document}
